% =====================================================================
% APPENDIX B: Addressing Complex Engineering Problems
% (RESTRUCTURE_PLAN_V2, appx-kp owner)
% Ported from old mainmatter/ch5.tex Section 5.2 (Addressing Complex
% Engineering Problems), with the K3/K4/K5/K8 knowledge-profile
% argument moved out to app_a.tex (\ref{app:kp}) and old Section 5.1's
% framing sentences folded into this appendix's opening paragraph.
% Table \ref{tab:cep_attributes} and its full P1-P7 content carried
% over intact; only the chapter cross-references were updated from the
% old hard-coded numbers to the new 9-chapter labels.
% =====================================================================
\chapter{Addressing Complex Engineering Problems}
\label{app:cep}

The Washington Accord's accreditation framework, adopted by the BAETE
program-outcome system, distinguishes routine engineering work from
work that addresses a Complex Engineering Problem (CEP). A problem is
complex when it cannot be resolved without in-depth engineering
knowledge. This complexity is captured by a set of P attributes, the
first of which, depth of knowledge, is mandatory. This appendix maps
each P attribute to the concrete parts of this thesis that exhibit it.

Designing the network at the center of this thesis is a complex
engineering problem in the formal sense: it cannot be solved by
natural science or mathematics alone. It required in-depth engineering
knowledge of the stereo pipeline, specialist knowledge of the
deep-stereo literature and its building blocks, the design of a novel
architecture under hard resource constraints, and grounding in a large
body of research literature to select and combine the right
mechanisms. Appendix~\ref{app:kp} maps these four demands onto the
Washington Accord's K3, K4, K5, and K8 knowledge profiles and shows
where we exercised each one.

Table~\ref{tab:cep_attributes} maps the P attributes of the Washington
Accord definition to the parts of this thesis that satisfy them.
Attribute P1 is mandatory and is carried by the knowledge profile
detailed in Appendix~\ref{app:kp}; the remaining rows record where the
thesis exhibits the other characteristics of a complex problem.

\begin{table}[!htbp]
\centering
\caption{Complex Engineering Problem attributes mapped to thesis sections.}
\label{tab:cep_attributes}
\small
\begin{tabular}{lp{4.2cm}p{7.2cm}}
\toprule
Attribute & Definition & Evidence in this thesis\\
\midrule
P1 & Depth of knowledge required & Cannot be solved without in-depth engineering knowledge (Appendix~\ref{app:kp}): K3 stereo fundamentals (Chapters~\ref{ch:lit}, \ref{ch:design}), K4 specialist deep-stereo knowledge (Chapter~\ref{ch:design}), K5 design under hard constraints (Chapters~\ref{ch:design}, \ref{ch:implementation}), K8 research-literature grounding (Chapter~\ref{ch:lit}).\\
P2 & Range of conflicting requirements & Accuracy, latency, memory, and parameter count pull against one another; the design criteria of Chapter~\ref{ch:design} demand balance across all of them jointly, and the studies of Chapter~\ref{ch:results} resolve the conflicts by measurement.\\
P3 & Depth of analysis required & No obvious correct architecture exists; we analyzed candidates through controlled ablation (Chapter~\ref{ch:results}) rather than selecting them by inspection, and rejected several plausible additions on evidence.\\
P4 & Familiarity of issues & The problem involves infrequently encountered issues: cross-domain generalization failure, surfaced through the test protocol of Chapter~\ref{ch:experimental} and documented in Chapter~\ref{ch:results}, and edge-inference optimization, both beyond routine practice (Chapter~\ref{ch:implementation}).\\
P5 & Extent of applicable codes & The work operates outside prescriptive engineering codes; community benchmark protocols and metric definitions (Chapter~\ref{ch:design}) are the closest standards, and we had to apply and match them carefully rather than follow them mechanically.\\
P6 & Extent of stakeholder involvement & The design serves stakeholders with differing needs: platform builders needing real-time latency, end users needing privacy and safety, and cost-sensitive deployments, needs examined in this thesis's socio-economic impact analysis.\\
P7 & Interdependence & The system is composed of interacting sub-parts, encoder, initialization, recurrent refinement, geometry volume, and upsampling, whose interactions the ablations of Chapter~\ref{ch:results} show are not separable.\\
\bottomrule
\end{tabular}
\end{table}

The attributes in the table are not asserted abstractly; each points
to a concrete artifact. The conflicting requirements of P2 are the
hard design envelope stated in Chapter~\ref{ch:design} and verified by
the efficiency measurements of Chapter~\ref{ch:results}. The depth of
analysis of P3 is the ablation record of Chapter~\ref{ch:results}, in
which we added candidate mechanisms one at a time under a matched
protocol and rejected several because they degraded accuracy. The
interdependence of P7 is visible in the same record: the value of one
block depends on which others are present, which is why we proceeded
by controlled composition rather than independent selection.

This appendix has shown that the thesis addresses a complex
engineering problem in the Washington Accord sense. The mandatory
attribute P1 is satisfied through the knowledge profile of
Appendix~\ref{app:kp}, and Table~\ref{tab:cep_attributes} grounds the
remaining P attributes in concrete parts of the thesis: the
conflicting requirements of the edge envelope, the ablation-driven
depth of analysis, and the interdependence of the architectural
blocks. Appendix~\ref{app:cea} turns to the activities we carried out
to reach this design and shows that they, too, meet the Washington
Accord's threshold for complexity.

% SOURCES: old mainmatter/ch5.tex Section 5.1 (Introduction, framing
% sentences folded into this appendix's opening two paragraphs),
% Section 5.2 (Addressing Complex Engineering Problems, CEP framing
% prose and Table tab:cep_attributes carried over intact with all
% seven P rows), and part of Section 5.4 (Summary, the P1/table
% grounding sentences folded into this appendix's closing paragraph).
% DROPPED: the K3/K4/K5/K8 detailed per-profile sentences that used to
% sit inside this section's opening paragraph; that material now lives
% in app_a.tex and is referenced here via \ref{app:kp} instead of
% repeated.
% MOVED-AWAY: none beyond what is logged above; the CEA material (old
% Section 5.3) is in app_c.tex, also owned by this agent. Old
% hard-coded "Chapter 2/3/7/8" numbers were replaced with \ref to the
% new chapter labels; the P6 row's old reference to the socio-economic
% chapter is phrased descriptively rather than as a numbered
% cross-reference because that chapter's label was not yet assigned
% when this appendix was written.
